
\documentclass{article}
\usepackage{amsmath, amsfonts, amssymb, amsthm} %packages for math-related
\usepackage[margin=1.2in]{geometry}
\usepackage[shortlabels]{enumitem}
\usepackage[utf8]{inputenc}
\usepackage{graphicx} % package for inserting images
\usepackage{commath} % package for things like \del, \cbr, and \sbr. These handle parentheses well.
\usepackage[mathscr]{euscript}%for \scr command 
\usepackage{../commands} %package with all of the commands for this class
\usepackage{url}
\setlength{\parindent}{0em} % so paragraphs aren't indented
\newcommand{\lcm}{\text{lcm}}
\newcommand{\Hom}{\text{Hom}}
\newcommand{\Ann}{\text{Ann}}
\newcommand{\Cc}{{C_c^\infty(\RR)}}
% ********************************************************** %
%               THINGS TO EDIT BELOW THIS LINE               %
% ********************************************************** %
\newcommand{\D}{\nabla}
\renewcommand{\d}{\partial}
\usepackage{wrapfig}
\title{\textsc{MATH 173 Problem Set 5}}
\author{Stepan (Styopa) Zharkov}
\date{May 4, 2022}
\begin{document}
\maketitle
\problem{1} Suppose that $f : \RR \to \CC$ is continuous with $(1 + |x|)^Nf(x)$ bounded for some $N > 1$ (or
indeed simply that $f \in L^1(\RR))$ and $a \in R$:
\begin{enumerate}[(a)]
    \item Let $f_a(x) = f(x - a)$. Show that $\hat{f}_a(y) = e^{-iay} \hat{f}(y)$.
\item Let $g_a(x) = e^{ixa}f(x)$. Show that $\hat{g}_a(y) = \hat{f}(y - a)$. 
\end{enumerate}
 \tri
\hop 
\solution
\begin{enumerate}[(a)]
    \item This problem is a simple computation. We see that, with a change of variables $z = x-a$, we have
    \begin{align*}
        \hat{f_a}(y) &= \cal{F}(f(x-a))(y) \\
        &= \int e^{-ixy}f(x-a)dx \\
        &= \int e^{-i(z+a)y}f(z)dx \\
        &= e^{iay}\int e^{-izy}f(z)dx \\
        &= e^{iay}\hat{f}(y),
    \end{align*}
    as we wanted. \qed
    \item This problem is even simpler computation. We see that 
    \begin{align*}
        \hat{g_a}(y) &= \cal{F}(e^{ixa}f(x))\\
        &= \int e^{-ixy}e^{ixa}f(x)dx \\
        &= \int e^{-ix(y-a)}f(x)dx \\
        &= \hat{f}(y-a),
    \end{align*}
    as we wanted. \qed
\end{enumerate}


\newpage
\problem{2} Compute the Fourier transform of the following functions:
\begin{enumerate}[(a)]
    \item $\chi_{(-a,a)}$
    \item $f(x) = e^{a|x|}, a > 0$
    \item $g(x) = |x|^ne^{-a|x|}, a>0, n \in \NN$. \tri
\end{enumerate} 
\solution
\begin{enumerate}[(a)]
    \item This problem is also computation.
    \begin{align*}
        (\cal{F}\chi_{(-a,a)}(y)) &= \int_{-\infty}^\infty e^{-ixy}\chi_{(-a,a)}(x) dx \\
        &= \int_{-a}^a e^{-ixy}dx \\
        &=\begin{cases}
            -\frac{i}{y}\del{e^{-iay}-e^{iay}} &\text{ if } y \ne 0\\
            2a &\text{ if } y =0
        \end{cases}
    \end{align*}
    \qed
    \item  Note that since $y \in \RR$ and $a > 0$, we know $iy - a \ne 0$ and $iy+a \ne 0$ so we can divide by them. So, 
    \begin{align*}
        (\cal{F}(e^{-a|x|})(y) &= \int_{-\infty}^{\infty}e^{-ixy}e^{a|x|}dx\\
        &= \int_{0}^{\infty}e^{-x(iy+a)}dx+\int_{-\infty}^{0}e^{-x(iy-a)}dx \\
        &= \sbr{-\frac{1}{iy+a}e^{-x(iy+a)}}_0^\infty+\sbr{-\frac{1}{iy-a}e^{-x(iy-a)}}_{-\infty}^0\\
        &= \frac{1}{iy+a} - \frac{1}{iy-a}
    \end{align*}
    because $a>0$. \qed
    \item In this problem, we will use repeated integration by parts. We see that executing integration by parts, we have
    \begin{align*}
        \cal{F}(|x|^ne^{-a|x|})(y) &= \int_{-\infty}^\infty |x|^ne^{-ixy-a|x|} dx\\
        &=  \int_{-\infty}^0 (-x)^ne^{-ixy+ax} dx+   \int_0^\infty x^ne^{-ixy-ax}dx\\
        &=  \int_{-\infty}^0 (-x)^ne^{-ix(y-a)} dx+   \int_0^\infty x^ne^{-ix(y+a)}dx \\
        &=  \frac{-n}{yi-a}\int_{-\infty}^0 (-x)^{n-1}e^{-ix(y-a)} dx+  \frac{n}{yi+a} \int_0^\infty x^{n-1}e^{-ix(y+a)}dx.
    \end{align*}
    Note that the boundary terms vanish in the integration by parts. Repeating integration by parts $n$ times, we see that
    \begin{align*}
        \cal{F}(|x|^ne^{-a|x|})(y) &= (-1)^n\frac{n!}{(yi-a)^n}\int_{-\infty}^0 e^{-ix(y-a)} dx+  \frac{n!}{(yi+a)^n} \int_0^\infty e^{-ix(y+a)}dx\\
        &=  (-1)^{n+1}\frac{n!}{(yi-a)^{n+1}}+  \frac{n!}{(yi+a)^{n+1}},
    \end{align*}
    and $yi\pm a$ does not vanish because $a>0$. \qed
\end{enumerate}



\newpage
\problem{3} 
\begin{enumerate}[(a)]
    \item Show that if $f$ is even (or odd) then so is its Fourier transform.
    \item Suppose that $f$ is even and real valued. Show that $\hat{f}$ is also real valued. What can you say about
    $\hat{f}$ when $f$ is odd and real valued.
\end{enumerate}
\tri
\hop 
\solution
\begin{enumerate}[(a)]
    \item The solution to this is straightforward. As Tadashi Tokieda would say, ``follow your nose''. First, let $f(x)= f(-x)$. Then, letting $z=-x$ power a change of variables, we see that 
    \begin{align*}
        \hat{f}(y) &= \int_{-\infty}^\infty e^{-ixy}f(x)dx\\
        &=  \int_{\infty}^{-\infty} -e^{izy}f(-z)dz\\
        &=  \int_{-\infty}^{\infty} e^{izy}f(-z)dz\\
        &=  \int_{-\infty}^{\infty} e^{-iz(-y)}f(z)dz\\
        &= \hat{f}(-y).
    \end{align*}
    Similarly, now let $f(x)= -f(-x)$. Then, 
    \begin{align*}
        \hat{f}(y) &= \int_{-\infty}^\infty e^{-ixy}f(x)dx\\
        &=  \int_{\infty}^{-\infty} -e^{izy}f(-z)dz\\
        &=  \int_{-\infty}^{\infty} e^{izy}f(-z)dz\\
        &=  \int_{-\infty}^{\infty}- e^{-iz(-y)}f(z)dz\\
        &= -\hat{f}(-y).
    \end{align*}
    So, the fourier transform preserves evenness and oddness. \qed
    \item Let $f$ be even. Then, 
    \begin{align*}
        \hat{f}(y) &= \int_{-\infty}^\infty e^{-ixy}f(x)dx \\
        &= \int_{-\infty}^0 e^{-ixy}f(x)dx + \int_{0}^\infty e^{-ixy}f(x)dx\\
        &= \int_{0}^\infty e^{ixy}f(-x)dx + \int_{0}^\infty e^{-ixy}f(x)dx\\
        &= \int_{0}^\infty e^{ixy}f(x)dx + \int_{0}^\infty e^{-ixy}f(x)dx\\
        &= \int_{0}^\infty\del{ e^{ixy}+ e^{-ixy}}f(x)dx\\
    \end{align*}
    This is real valued because $ e^{ixy}+ e^{-ixy} \in \RR$ and $f(x) \in \RR$. 
    \hop 
    The idea is the same when $f$ is odd. By a similar computation,
    \begin{align*}
        \hat{f}(y) &= \int_{-\infty}^\infty e^{-ixy}f(x)dx \\
        &= \int_{-\infty}^0 e^{-ixy}f(x)dx + \int_{0}^\infty e^{-ixy}f(x)dx\\
        &= \int_{0}^\infty e^{ixy}f(-x)dx + \int_{0}^\infty e^{-ixy}f(x)dx\\
        &= \int_{0}^\infty -e^{ixy}f(x)dx + \int_{0}^\infty e^{-ixy}f(x)dx\\
        &= \int_{0}^\infty\del{ -e^{ixy}+ e^{-ixy}}f(x)dx\\
    \end{align*}
    We see that this integral has no real part because $-e^{ixy}+ e^{-ixy}$ has no real part and $f(x)\in \RR$. \qed
\end{enumerate}


\newpage
\problem{4} Define the Hermite functions by $h_n = (x - d
dx
)^n e^{-x^2/2}$, such that $h_0 = e^{-x^2/2}$ and
$h_{n+1}(x) = xh_n(x) - h'_n(x)$. Show that $\hat{h}_n = (-i)^n\sqrt{2\pi} h_n$. \tri
\hop 
\solution
This proof follows by induction. From class we know that 
\[\hat{h_0}(y)=\sqrt{2\pi}e^{-y^2/2} = \sqrt{2\pi}h_0(y).\] 
So, the base case holds. Now, suppose the statement is true up to $h_n$. Then, we see that by linearity of $\cal{F}$ and the properties we have seen in class,
\begin{align*}
    \cal{F}(h_{n+1}) &= \cal{F}(xh_n)-\cal{F}(h_n')\\
    &= i(\cal{F}h_n)' - iy(\cal{F}h_n) \\
    &= -i(y\sqrt{2\pi}(-i)^n h_n - (\sqrt{2\pi}(-i)^n h_n)') \\
    &= (-i)^{n+1}\sqrt{2\pi}(yh_n - h_n')\\
    &= (-i)^{n+1}\sqrt{2\pi}h_{n+1}.
\end{align*}
So, by induction, the statement holds for all $n$. \qed

\newpage
\problem{5} Prove the Poisson summation formula
\[\sum_
{n\in\ZZ}
f(n)e^{-i\pi n} = \sum_
{m\in \ZZ}\hat{f}(x + 2\pi m),\]
where $f$ satisfies $|f(x)|+ |\hat{f}(x)|\le C(1 + |x|)^{-a}$ for some $a > 1$.  \tri
\hop 
\solution
We will follow the hint. Let us compute the Fourier coefficients of the RHS. We get that with a change of variables,
\begin{align*}
    c_n  &= \frac{1}{2\pi} \int_{-\pi}^\pi e^{inx}\sum_{m \in \ZZ}\hat{f}(x+2\pi m) dx \\
    &= \frac{1}{2\pi} \sum_{m \in \ZZ} \int_{-\pi}^\pi e^{inx}\hat{f}(x+2\pi m) dx \\
    &= \frac{1}{2\pi} \sum_{m \in \ZZ} \int_{-\pi+2\pi m}^{\pi+2\pi m} e^{in(z - 2\pi m)}\hat{f}(z) dz \\
    &= \frac{1}{2\pi} \sum_{m \in \ZZ} \int_{-\pi+2\pi m}^{\pi+2\pi m} e^{inz}\hat{f}(z) dz \\
    &= \frac{1}{2\pi} \int e^{inz}\hat{f}(z) dz \\
    &= \cal{F}\inv(\hat{f})(n) \\
    &= f(n).
\end{align*}
Here we also used that the Fourier transform reverts the action of the Fourier transform given these conditions, as given in the notes. So, assuming that the Fourier series converges to the function, 
\begin{align*}
    \sum_{m \in \ZZ}\hat{f}(x+2\pi m) &=  \sum_{n \in \ZZ}c_ne^{-inx}\\
    &= \sum_{n \in \ZZ}f(n)e^{-inx},
\end{align*}
which is what we wanted to prove. \qed
\begin{center}
    \includegraphics[height=3cm]{../images/fourier-people}
\end{center}

\newpage
\problem{6} Prove the Whittaker–Shannon interpolation formula
\[f(x) = \sum_{n\in \ZZ}
f(n)\text{sinc}(\pi(x -n))\]
provided $\supp(\hat{f}) \sse [-\pi,\pi]$ and $|f(x)|\le C(1 + |x|)^{-a}$. Here $\text{sinc}(x) = \sin(x)/x$.  \tri
\hop 
\solution
Since $\supp(\hat{f}) \sse [-\pi, \pi]$, we see that for $x \in (-\pi,\pi)$, by problem 5,
\begin{align*}
    \sum_{n \in \ZZ}f(n) e^{-inx} = \sum_{m \in \ZZ} \hat{f}(x + 2\pi m) = \hat{f}(x).
\end{align*}
Note that we saw that a finite number of points do not change the fourier transform, so the points $-\pi, \pi$ are not a problem. Assuming that the inverse of the Fourier transform reverts the action of the Fourier transform, as in problem 5, we see that 
\begin{align*}
    f(y) &= \cal{F}\inv(\hat{f})(y) \\
    &= \frac{1}{2\pi}\int e^{inx}\sum_{n\in \ZZ} f(n)e^{-ixy}dx \\
    &= \frac{1}{2\pi}\int_{-\pi}^\pi e^{ixy}\sum_{n\in \ZZ} f(n)e^{-inx}dx \\
    &=\sum_{n\in \ZZ} f(n) \frac{1}{2\pi}\int_{-\pi}^\pi e^{ixy}e^{-inx}dx \\
    &=\sum_{n\in \ZZ} f(n) \frac{1}{2\pi}\int_{-\pi}^\pi e^{ix(y-n)}dx \\
    &=\sum_{n\in \ZZ} f(n) \frac{1}{2\pi}\sbr{\frac{1}{i(y-n)}e^{ix(y-n)}}_{-\pi}^\pi dx \\
    &=\sum_{n\in \ZZ} f(n) \del{\frac{e^{i\pi(y-n)}-e^{-i\pi(y-n)}}{2\pi i(y-n)}} dx\\
    &=\sum_{n\in \ZZ} f(n) \del{\frac{\sin(\pi(y-n))}{\pi (y-n)}} dx\\
    &=\sum_{n\in \ZZ} f(n) \text{sinc}(\pi(y-n)),\\
\end{align*}
exactly as we wanted.\qed



\newpage
\problem{7} 
\begin{enumerate}[(a)]
    \item Use the Fourier transform to solve the following free Schr\" odinger equation $iu_t = -u_{xx}$,
    with the initial condition $u(0,x) = g(x)$, where $u : \RR\times \RR \to \CC$ and $g \in \cal{S}(\RR)$ is a given Schwarz function.
    You may leave your solution as the inverse Fourier transform of a Schwartz function (you do not need
    to evaluate it explicitly).
    \item  Evaluate the solution in part (a) explicitly if $g(x) = \frac{1}{\sqrt{2\pi}} e^{-x^2/
    2}$, and show that 
    \[\int_\RR|u(x,t)|^2dx = 1\] for
    all $t$.
\end{enumerate}
\tri
\hop 
\solution
\begin{enumerate}[(a)]
    \item Let $s = it$. Then, our equation becomes 
    \[iu_t = u_s = -u_{xx}.\]
    This is the heat equation, which we already know how to solve from class. We know 
    \begin{align*}
        u(t,x) = \del{\frac{1}{\sqrt{4\pi s}}e^{-\frac{z^2}{4s}}* g(z)}(x) =\del{\frac{1}{\sqrt{4\pi i t}}e^{-\frac{z^2}{4it}}* g(z)}(x).
    \end{align*}
    This is the answer we are looking for.\qed
    \item This problem is simply bashing out the convolution. Plugging in $g(z) = \frac{1}{\pi^{1/4}}e^{-\frac{z^2}{2}}$, we get that 
    \begin{align*}
        u(t,x) &= \del{\frac{1}{\sqrt{4\pi i t}}e^{-\frac{z^2}{4it}}* \frac{1}{\pi^{1/4}}e^{-\frac{z^2}{2}}}(x) \\
        &= \frac{1}{\pi^{1/4}\sqrt{4\pi i t}} \del{e^{-\frac{z^2}{4it}}* e^{-\frac{z^2}{2}}}(x) \\
        &=  \frac{1}{\pi^{1/4}\sqrt{4\pi i t}} \int \exp\del{ -\frac{(x-y)^2}{4ti} - \frac{y^2}{2}}dy\\
        &=  \frac{1}{\pi^{1/4}\sqrt{4\pi i t}} \int \exp\del{ -\frac{x^2-2xy+y^2+2tiy^2}{4ti}}dy\\
        &=  \frac{1}{\pi^{1/4}\sqrt{4\pi i t}} \int \exp\del{ -\frac{\del{\sqrt{2ti+1}y + \frac{1}{\sqrt{2ti+1}x}}^2}{4ti}}\exp\del{\frac{itx^2}{4t^2+1} - \frac{x^2}{8t^2+2}}dy\\
        &=  \frac{1}{\pi^{1/4}\sqrt{4\pi i t}}\exp\del{\frac{itx^2}{4t^2+1} - \frac{x^2}{8t^2+2}} \int \exp\del{ -\frac{\del{\sqrt{2ti+1}y + \frac{1}{\sqrt{2ti+1}x}}^2}{4ti}}dy\\
        &= \frac{1}{i\pi^{1/4}}\sqrt{\frac{1-2it}{1+4t^2}} \exp\del{\frac{itx^2}{4t^2+1} - \frac{x^2}{8t^2+2}}.
    \end{align*}
    Note that we solved this by completing the square to bring it to a form where we can use the Gaussian integral. 
    \hop
    For the second part of the problem, we will again use the Gaussian integral. Note that $e^{is}$ has absolute value 1 for real $s$, so we see that 
    \begin{align*}
        \int \abs{u(t,x)}^2 dx&=  \int \abs{\frac{1}{i\pi^{1/4}}\sqrt{\frac{1-2it}{1+4t^2}} \exp\del{\frac{itx^2}{4t^2+1} - \frac{x^2}{8t^2+2}}}^2dx \\
        &= \frac{1}{\pi^{1/2}}\abs{\frac{1-2it}{1+4t^2}}\int \abs{ \exp\del{\frac{itx^2}{4t^2+1} - \frac{x^2}{8t^2+2}}}^2dx \\
        &= \frac{1}{\pi^{1/2}}\abs{\frac{1-2it}{1+4t^2}}\int \abs{ \exp\del{\frac{x^2}{8t^2+2}}}^2dx \\
        &= \frac{1}{\pi^{1/2}}\abs{\frac{1-2it}{1+4t^2}}\int \exp\del{\frac{x^2}{4t^2+1}}dx \\
        &= \frac{1}{\pi^{1/2}}\frac{\sqrt{1+4t^2}}{1+4t^2}\int \exp\del{\frac{x^2}{4t^2+1}}dx\\
        &= \frac{1}{\pi^{1/2}}\frac{\sqrt{1+4t^2}}{1+4t^2}\sqrt{\pi (4t^2 + 1)}\\
        &= 1,
    \end{align*}
    exactly as we wanted. \qed
\end{enumerate}
\end{document}